% !PRINTED 472
\weberpage{472}

Auf Grund der Eigenschaften 1., 2., 3.\ nennen wir die Gesammtheit der $\mu$ Substitutionen des Normalkörpers $\Omega(\varrho)$ eine \gesperrt{Gruppe}, und zwar eine \gesperrt{endliche Gruppe} vom \gesperrt{Grade} $\mu$. Sie heisst auch die \gesperrt{Gruppe der Substitutionen des Körpers} $\Omega(\varrho)$.

Mit der allgemeinen Theorie der Gruppen werden wir uns in späteren Abschnitten eingehend beschäftigen.

\begin{center}
\S.\ 148.
\end{center}

\begin{center}
\gesperrt{Permutationsgruppen.}
\end{center}

Im \S.\ 145 haben wir gesehen, wie man zu einer beliebigen Gleichung in $\Omega$ vom Grade $m$
\weq{1}{F(x) = 0,}
von der nur vorausgesetzt wird, dass sie keine gleichen Wurzeln hat, einen Normalkörper $\Omega(\varrho)$ bestimmen kann, dem alle Wurzeln von (1) angehören. Wenn also die Wurzeln von (1) mit
\weq{2}{\alpha, \alpha_1, \alpha_2 \ldots \alpha_{m-1}}
bezeichnet werden, so ist
\weq{3}{\alpha = \chi(\varrho), \quad \alpha_1 = \chi_1(\varrho), \ldots \alpha_{m-1} = \chi_{m-1}(\varrho),}
und die $\chi$ sind rationale Functionen in $\Omega$.

Setzen wir aber für $\varrho$ eine mit $\varrho$ conjugirte Grösse $\varrho_1$, führen wir also eine Substitution $\sigma = (\varrho, \varrho_1)$ aus, so geht jede Wurzel $\alpha$ in eine bestimmte andere über, und niemals zwei verschiedene in dieselbe; denn wenn
\[
F(\alpha) = F[\chi(\varrho)] = 0
\]
ist, so ist nach \S.\ 146, 1:
\[
F[\chi(\varrho_1)] = 0,
\]
d.\,h. $\chi(\varrho_1)$ ist auch eine Wurzel der Gleichung (1). Dass aber nicht zwei verschiedene Grössen des Körpers $\Omega(\varrho)$ durch eine Substitution $\sigma$ in die gleiche Grösse übergehen können, haben wir im \S.\ 146 schon allgemein gezeigt.

\gesperrt{1. Eine Substitution $\sigma$ ruft also unter den Wurzeln $\alpha$ eine gewisse Permutation hervor, indem jeder der Indices $0, 1, 2 \ldots m-1$ durch einen bestimmten anderen ersetzt wird.}
